\section*{Capítulo \ref{ch:b1:carbohydrates} --- Glúcidos}

\begin{solution}{exo:b1:carbohydrates:1}
Glucosa: hexosa, \omterm{def:b1:carbohydrates:monosaccharide}{aldosa}. Fructosa: hexosa, \omterm{def:b1:carbohydrates:monosaccharide}{cetosa}. Ribosa: pentosa, \omterm{def:b1:carbohydrates:monosaccharide}{aldosa}.
Gliceraldehído: triosa, \omterm{def:b1:carbohydrates:monosaccharide}{aldosa}.
\end{solution}

\begin{solution}{exo:b1:carbohydrates:2}
El carbono carbonílico tras el cierre del anillo, que lleva el oxígeno del
anillo y un hidroxilo y puede volver a abrirse a aldehído o a cetona. En la
sacarosa los dos carbonos anoméricos están bloqueados en el \omterm{def:b1:carbohydrates:glycosidic}{enlace
glucosídico}; ninguno puede abrirse, de modo que no se forma aldehído y el
cobre no se reduce.
\end{solution}

\begin{solution}{exo:b1:carbohydrates:3}
\omterm{def:b1:carbohydrates:polysaccharide}{Almidón}: glucosa, $\alpha(1{\to}4)$ con ramificaciones $\alpha(1{\to}6)$
(amilopectina). \omterm{def:b1:carbohydrates:polysaccharide}{Glucógeno}: lo mismo, más ramificado. \omterm{def:b1:carbohydrates:polysaccharide}{Celulosa}: glucosa,
$\beta(1{\to}4)$. \omterm{def:b1:carbohydrates:polysaccharide}{Quitina}: N-acetilglucosamina, $\beta(1{\to}4)$.
\end{solution}

\begin{solution}{exo:b1:carbohydrates:4}
C1 abajo, C2 abajo, C3 arriba, C4 abajo. La $\alpha$-D-galactosa es igual
con el hidroxilo del C4 arriba.
\end{solution}

\begin{solution}{exo:b1:carbohydrates:5}
En el enlace $\alpha$ el hidroxilo anomérico apunta hacia abajo, y cada
residuo puede añadirse con el mismo ángulo que el anterior: la cadena gira
de manera constante y se cierra en hélice. En el enlace $\beta$ apunta hacia
arriba, y un residuo solo puede añadirse volteado: la cadena queda recta.
Las moléculas de yodo caben dentro de la hélice de amilosa y su estado
electrónico cambia (azul); una cinta plana no ofrece ninguna cavidad.
\end{solution}

\begin{solution}{exo:b1:carbohydrates:6}
Yodo: solo el \omterm{def:b1:carbohydrates:polysaccharide}{almidón} vira a azul negruzco. Benedict sobre los otros dos:
solo la glucosa da un precipitado rojo; la sacarosa sigue azul. Confírmese
la sacarosa calentándola con ácido diluido, neutralizando y repitiendo
Benedict: ahora es positivo.
\end{solution}

\begin{solution}{exo:b1:carbohydrates:7}
Cadenas de 12 residuos: unas \num{4000} cadenas, la mitad de ellas
externas: unos \num{2000} extremos. Amilosa: un extremo. Para liberar
\num{5000} glucosas a una por segundo y por extremo: el \omterm{def:b1:carbohydrates:polysaccharide}{glucógeno},
\qty{2.5}{s}; la amilosa, \qty{5000}{s}, casi hora y media.
\end{solution}

\begin{solution}{exo:b1:carbohydrates:8}
$\sigma = 0.8\times 25/(2\times 0.5) = \qty{20}{MPa}$: la cincuentava parte
de la resistencia de la fibrilla; la pared tiene un amplio margen de
seguridad, necesario porque las fibrillas son solo la cuarta parte de la
pared.
\end{solution}

\begin{solution}{exo:b1:carbohydrates:9}
La lactasa rompe el enlace $\beta(1{\to}4)$ entre la galactosa y la glucosa;
sin ella la lactosa no se absorbe (solo los \omterm{def:b1:carbohydrates:monosaccharide}{monosacáridos} atraviesan el
\omterm{def:b1:body-plans-tissues:epithelium}{epitelio}). En el colon las bacterias la fermentan a ácidos y gas
(retortijones, distensión), y el azúcar sin digerir y los ácidos atraen agua
por \omterm{def:b1:membranes-transport:osmosis}{ósmosis} hacia la luz: diarrea.
\end{solution}

\begin{solution}{exo:b1:carbohydrates:10}
Las bacterias y los protistas del rumen secretan celulasas y fermentan la
glucosa a \omterm{def:b1:lipids:fattyacid}{ácidos grasos} cortos, que la vaca absorbe y oxida; los propios
microbios se digieren después. La fermentación debe preceder al intestino
delgado para que sus productos lleguen a la superficie absorbente y la
proteína microbiana llegue al estómago y al intestino. \omterm{def:b1:carbohydrates:polysaccharide}{Celulosa}: \qty{4}{kg};
fermentada, \qty{2.4}{kg}; \qty{29}{MJ} al día, casi toda la energía de la
vaca.
\end{solution}

\begin{solution}{exo:b1:carbohydrates:11}
Las \omterm{def:b1:cell-unit-of-life:cell}{células} O no llevan ningún antígeno y no provocan anticuerpos en nadie;
el \omterm{def:b1:mammal-organization:compartments}{plasma} O, en cambio, contiene anticuerpos contra A y contra B, de modo
que una persona O ataca cualquier \omterm{def:b1:cell-unit-of-life:cell}{célula} A o B que reciba. El sistema
inmunitario de una persona A nunca ha visto B y fabrica anticuerpos anti-B
(contra bacterias intestinales que llevan azúcares parecidos), pero tolera
A: el \omterm{def:b1:mammal-organization:compartments}{plasma} A contiene solo anti-B.
\end{solution}

\begin{solution}{exo:b1:carbohydrates:12}
Acierta en que la glucosa es el combustible de casi toda \omterm{def:b1:cell-unit-of-life:cell}{célula} y la
\omterm{def:b1:carbohydrates:polysaccharide}{celulosa} la molécula orgánica más abundante de la Tierra, y en que ambas son
polímeros de D-glucosa, o la propia D-glucosa. Un solo enlace lo cambia
todo: $\alpha$ da un ovillo que las enzimas abren y las \omterm{def:b1:cell-unit-of-life:cell}{células} queman;
$\beta$ da una fibra que resiste las enzimas, el agua y la tracción, y que
la mayoría de los animales no pueden aprovechar en absoluto. La biología no
es la lista de las moléculas, sino la geometría de su unión: la misma
química puede ser alimento o andamio según una única elección
estereoquímica, y a los \omterm{def:b1:organism-environment:organism}{organismos} los separa qué enzimas poseen para
deshacerla.
\end{solution}

\begin{solution}{pb:b1:carbohydrates:1}
\textbf{1.} $2^{t-1}$: 1, 2, 4 y $2^{11} = 2048$.
\textbf{2.} $2^{12} - 1 = 4095$.
\textbf{3.} $4095\times 13 = \num{53235}$, unos \num{53000}.
\textbf{4.} $\num{53235}\times 162 = \qty{8.6e6}{Da}$;
$\qty{1.43e-17}{g}$.
\textbf{5.} $2048/4095$: la mitad. \num{2048} extremos no reductores.
\textbf{6.} Uno.
\textbf{7.} $12\times 1.9 = \qty{23}{nm}$ de radio, \qty{46}{nm} de
diámetro: el doble que un ribosoma.
\textbf{8.} $100/\num{1.43e-17} = \num{7.0e18}$ partículas.
\textbf{9.} $100/162 = \qty{0.62}{mol}$ de residuos, es decir, \qty{111}{g}
de glucosa libre (el agua de la \omterm{prop:b1:water-small-molecules:families}{hidrólisis} aporta la diferencia).
\textbf{10.} \qty{400}{g}, más de la cuarta parte de la masa del hígado.
\textbf{11.} \qty{10}{g/h} $= \qty{2.78}{mg/s} = \qty{1.54e-5}{mol/s} =
\num{9.3e18}$ moléculas por segundo.
\textbf{12.} Unas \qty{11}{h} (111 g a 10 g/h): una noche. Sobre todo al
encéfalo, que consume \qty{5}{g/h}.
\textbf{13.} $2048\times\num{7e18} = \num{1.4e22}$ extremos.
\textbf{14.} $\num{1.4e22}\times 20 = \num{2.9e23}$ por segundo $=
\qty{0.48}{mol/s} = \qty{86}{g/s}$: treinta mil veces la velocidad real. Los
extremos nunca son limitantes; la enzima sí.
\textbf{15.} $\num{9.3e18}/20 = \num{4.6e17}$ moléculas, unos
\qty{75}{mg} de enzima, una fracción pequeña de la proteína del hígado.
\textbf{16.} Cadenas de \num{53000} residuos: siguen siendo $\num{7e18}$
extremos (uno cada una), \num{2048} veces menos; velocidad máxima
\qty{42}{mg/s}, todavía quince veces la real. (Con cadenas de amilosa más
largas el margen se estrecha; el límite real de una reserva sin ramificar es
su insolubilidad y su cristalización, que dejan los extremos fuera de
alcance.)
\textbf{17.} Cuando la glucemia cae, las hormonas activan la fosforilasa en
segundos; la enzima encuentra miles de extremos por partícula ya expuestos
en la superficie, de modo que la velocidad de liberación puede multiplicarse
por cien de golpe, sin esperar a que se desenrollen cadenas ni a que se
disuelvan partículas.
\textbf{18.} $0.55/1.0 = \qty{0.55}{osmol/L}$ (\qty{100}{g}
$= \qty{0.55}{mol}$ de glucosa libre).
\textbf{19.} Casi el doble de la osmolaridad del \omterm{def:b1:eukaryotic-cell:organelle}{citosol}: el agua entraría a
raudales desde la sangre, las \omterm{def:b1:cell-unit-of-life:cell}{células} se hincharían hasta casi el triple de
su volumen y reventarían.
\textbf{20.} $\num{7e18}/\num{6e23} = \qty{1.2e-5}{mol}$ en \qty{1}{L}:
\qty{12}{\micro osmol/L}, cuarenta mil veces menos.
\textbf{21.} $\Psi_s = -2.58\times 0.55 = \qty{-1.4}{MPa}$: una pared
tendría que soportar \qty{1.4}{MPa}, catorce atmósferas.
\textbf{22.} El \omterm{def:b1:carbohydrates:polysaccharide}{glucógeno} muscular es el combustible propio del músculo para
la contracción, movilizado como glucosa-6-fosfato directamente a la
glucólisis; el músculo carece de la enzima que libera la glucosa de su
fosfato, de modo que la reserva no puede salir de la fibra.
\textbf{23.} Una patata moviliza su \omterm{def:b1:carbohydrates:polysaccharide}{almidón} a lo largo de semanas de
brotación; un hígado debe responder en minutos. Los granos grandes y densos
con pocos extremos en la superficie convienen a una liberación lenta y
sostenida; las partículas pequeñas y muy ramificadas con miles de extremos,
a una rápida.
\textbf{24.} $0.62/(4\times 3600) = \qty{4.3e-5}{mol/s} = \num{2.6e19}$
residuos por segundo; \qty{1.24}{mol} de \omterm{def:b1:water-small-molecules:atp}{ATP} en total, $\num{5.2e19}$ por
segundo.
\textbf{25.} Liberación: $\num{9.3e18}$ moléculas de glucosa por segundo
(\qty{10}{g/h}); la hacen posible \num{1.4e22} extremos no reductores,
\num{2048} por partícula; y el polímero ahorra a las \omterm{def:b1:cell-unit-of-life:cell}{células}
\qty{0.55}{osmol/L}, casi el doble de su propia osmolaridad.
\end{solution}
